% Provenance (never printed): NSI premiere, "Representation des donnees: types construits" (p-uplets, p-uplets
% nommes, tableau donne en comprehension); NSI terminale "Mise au point" (side effects through aliasing).
% See sources/book2/BRIEFS.md.
\chapter{Tuples, Sets and Comprehensions}\label{ch:g10:tuples-sets-comprehensions}

Dara and Umar both borrow books from the town library, and they wonder whether they have the same taste. The
library's records give each of them a list of loans; a book borrowed twice appears twice. To answer ``which books
have we both read?'' with lists, one has to take each book of Dara's list and look for it along the whole of Umar's,
and then remove the repeats from the answer.

Python has a type made for this question, the \emph{\omterm{def:g10:tuples-sets-comprehensions:set}{set}}, where a book is either in or out and never twice. This
chapter introduces it, together with the \emph{\omterm{def:g10:tuples-sets-comprehensions:tuple}{tuple}}, a fixed group of values such as a point $(3, 4)$, and the
\emph{\omterm{def:g10:tuples-sets-comprehensions:comprehension}{comprehension}}, a one-line way to build a list, a \omterm{def:g10:tuples-sets-comprehensions:set}{set} or a dictionary. On the way it answers a question that
caused a bug in the previous chapter: why does a function that sorts its argument change the caller's list?

\begin{recall}
A list holds values in order, \code{t[i]} reads the value at position \code{i}, and \code{t.append(x)} adds
\code{x} at the end (Book 1, grade 7). A dictionary maps keys to values, \code{d[k]}, and a record of a table is a
dictionary from column names to values (Book 1, grade 8). The town library's tables were met in Book 1; they are
read here with \code{town.load("loans")}, which returns a list of such records. \omterm{def:g10:tuples-sets-comprehensions:set}{Sets} as mathematical objects belong
to the mathematics book; this chapter is about the Python type.
\end{recall}

\section{Tuples}\label{sec:g10:tuples-sets-comprehensions:tuples}

\begin{definition}[Tuple, unpacking]\label{def:g10:tuples-sets-comprehensions:tuple}
A \emph{tuple}\index{tuple} is a sequence of values written between parentheses and separated by commas, such as
\code{(3, 4)}. Like a list it has a length and positions, read with \code{p[i]}; unlike a list, it cannot be changed
once built. \emph{Unpacking}\index{unpacking} assigns the values of a tuple to as many names at once:
\code{x, y = p}.
\end{definition}

\omcode{code/grade-10/02-tuples-sets-comprehensions/listings/tuples.py}{1}{14}{Building, reading and \omterm{def:g10:tuples-sets-comprehensions:tuple}{unpacking}
\omterm{def:g10:tuples-sets-comprehensions:tuple}{tuples}.}\label{lst:g10:tuples-sets-comprehensions:tuples}

\omoutput{code/grade-10/02-tuples-sets-comprehensions/outputs/tuples.txt}{A \omterm{def:g10:tuples-sets-comprehensions:tuple}{tuple} of one element needs a comma;
without it the parentheses are only grouping.}\label{out:g10:tuples-sets-comprehensions:tuples}

Line 9 is the classic use of \omterm{def:g10:tuples-sets-comprehensions:tuple}{unpacking}: the right side \code{b, a} builds the \omterm{def:g10:tuples-sets-comprehensions:tuple}{tuple} \code{(2, 1)} first, and only
then are its two values given to \code{a} and \code{b}. No third variable is needed to swap two values.

The most common use of a \omterm{def:g10:tuples-sets-comprehensions:tuple}{tuple} is to return several results from one function. Strictly, a function still returns
one value, a \omterm{def:g10:tuples-sets-comprehensions:tuple}{tuple}; the caller unpacks it.

\omcode{code/grade-10/02-tuples-sets-comprehensions/listings/min_max.py}{1}{15}{One loop, two results, returned as
a \omterm{def:g10:tuples-sets-comprehensions:tuple}{tuple}.}\label{lst:g10:tuples-sets-comprehensions:min-max}

\omoutput{code/grade-10/02-tuples-sets-comprehensions/outputs/min_max.txt}{The second call prints the \omterm{def:g10:tuples-sets-comprehensions:tuple}{tuple}
itself.}\label{out:g10:tuples-sets-comprehensions:min-max}

A \omterm{def:g10:tuples-sets-comprehensions:tuple}{tuple} whose positions have fixed meanings (a title, an author, a year) is a record. Numbers are easy to mix up;
names are not.

\begin{definition}[Named tuple]\label{def:g10:tuples-sets-comprehensions:named-tuple}
A \emph{named tuple}\index{named tuple} is a \omterm{def:g10:tuples-sets-comprehensions:tuple}{tuple} whose positions also have names, so that \code{b.year} and
\code{b[2]} read the same value. In Python, \code{namedtuple("Book", ["title", "author", "year"])}, from the module
\code{collections}, creates a new type \code{Book} of such \omterm{def:g10:tuples-sets-comprehensions:tuple}{tuples}.
\end{definition}

\omcode{code/grade-10/02-tuples-sets-comprehensions/listings/named.py}{1}{10}{A record of the town library as a
\omterm{def:g10:tuples-sets-comprehensions:named-tuple}{named tuple}.}\label{lst:g10:tuples-sets-comprehensions:named}

\omoutput{code/grade-10/02-tuples-sets-comprehensions/outputs/named.txt}{A \omterm{def:g10:tuples-sets-comprehensions:named-tuple}{named tuple} prints its field names, and
still unpacks like any \omterm{def:g10:tuples-sets-comprehensions:tuple}{tuple}.}\label{out:g10:tuples-sets-comprehensions:named}

A dictionary could hold the same record. The \omterm{def:g10:tuples-sets-comprehensions:named-tuple}{named tuple} is fixed: it always has the same three fields and they
cannot be changed, which is what a record read from a table usually needs.

\section{Mutable and immutable values}\label{sec:g10:tuples-sets-comprehensions:mutable}

\begin{definition}[Mutable, immutable]\label{def:g10:tuples-sets-comprehensions:mutable}
A value is \emph{mutable}\index{mutable} when it can be changed after it is built, and
\emph{immutable}\index{immutable} when it cannot. Lists, dictionaries and \omterm{def:g10:tuples-sets-comprehensions:set}{sets} are mutable. Numbers, strings and
\omterm{def:g10:tuples-sets-comprehensions:tuple}{tuples} are immutable: an operation on them builds a new value.
\end{definition}

\omcode{code/grade-10/02-tuples-sets-comprehensions/listings/immutable.py}{1}{7}{The same assignment on a list and
on a \omterm{def:g10:tuples-sets-comprehensions:tuple}{tuple}.}\label{lst:g10:tuples-sets-comprehensions:immutable}

\omoutput{code/grade-10/02-tuples-sets-comprehensions/outputs/immutable.txt}{Changing a position of a
\omterm{def:g10:tuples-sets-comprehensions:tuple}{tuple} is refused.}\label{out:g10:tuples-sets-comprehensions:immutable}

The difference matters because of the way Python names values. A variable does not contain a list; it
\emph{refers} to a list stored somewhere in memory. The assignment \code{b = a} does not copy anything: it makes
\code{b} refer to the same list as \code{a}.

\begin{definition}[Alias]\label{def:g10:tuples-sets-comprehensions:alias}
Two names are \emph{aliases}\index{alias} when they refer to the same value in memory. The test \code{a is b} is
true exactly when \code{a} and \code{b} are aliases; the test \code{a == b} compares the contents of two values,
which may be distinct.
\end{definition}

\omcode{code/grade-10/02-tuples-sets-comprehensions/listings/alias.py}{1}{9}{An \omterm{def:g10:tuples-sets-comprehensions:alias}{alias} and a
copy.}\label{lst:g10:tuples-sets-comprehensions:alias}

\omoutput{code/grade-10/02-tuples-sets-comprehensions/outputs/alias.txt}{A change made through \code{b} is seen
through \code{a}; a change made to the copy \code{c} is not.}\label{out:g10:tuples-sets-comprehensions:alias}

The figure below draws the memory after line 7: names on the left, values on the right, an arrow for ``refers
to''.

\begin{omfigure}
\begin{tikzpicture}[font=\small, cell/.style={draw, minimum width=7mm, minimum height=6mm, inner sep=0pt,
    font=\small\ttfamily}, nm/.style={draw, rounded corners=2pt, fill=omDef!8, minimum width=8mm,
    minimum height=6mm, font=\small\ttfamily}]
  \node[nm] (a) at (0,1.2) {a};
  \node[nm] (b) at (0,0.2) {b};
  \node[nm] (c) at (0,-1.4) {c};
  \foreach \v [count=\i] in {1,2,3,4} \node[cell] (L\i) at (3+0.7*\i,0.7) {\v};
  \foreach \v [count=\i] in {1,2,3,4,5} \node[cell, fill=omProp!10] (M\i) at (3+0.7*\i,-1.4) {\v};
  \draw[ptr] (a.east) -- (L1.west);
  \draw[ptr] (b.east) -- (L1.west);
  \draw[ptr] (c.east) -- (M1.west);
  \node[anchor=west, font=\footnotesize, align=left] at (6.7,0.7) {one list, two names:\\\code{a is b} is \code{True}};
  \node[anchor=west, font=\footnotesize, align=left] at (7.4,-1.4) {a second list, built by\\\code{list(a)}};
\end{tikzpicture}

{\small Names refer to values. \code{a} and \code{b} are \omterm{def:g10:tuples-sets-comprehensions:alias}{aliases} of one list; \code{c} refers to a copy, which
changes on its own.}
\end{omfigure}

This explains the bug of the previous chapter: inside \code{largest_three(t)}, the parameter \code{t} is an \omterm{def:g10:tuples-sets-comprehensions:alias}{alias}
of the caller's list, so \code{t.sort()} sorts the caller's list. \omterm{def:g10:tuples-sets-comprehensions:alias}{Aliases} of an \omterm{def:g10:tuples-sets-comprehensions:mutable}{immutable} value are harmless,
since nobody can change it; \omterm{def:g10:tuples-sets-comprehensions:alias}{aliases} of a \omterm{def:g10:tuples-sets-comprehensions:mutable}{mutable} value are where \omterm{def:g10:specifications-and-tests:side-effect}{side effects} come from.

A copy made by \code{list(a)}, or by the slice \code{a[:]}, is a new list holding the same elements. When those
elements are themselves lists, they are not copied: the new outer list refers to the same inner lists.

\begin{definition}[Shallow copy, deep copy]\label{def:g10:tuples-sets-comprehensions:copy}
A \emph{shallow copy}\index{shallow copy} of a value is a new value whose parts refer to the same objects as the
parts of the original. A \emph{deep copy}\index{deep copy} copies the parts too, and their parts, all the way down,
so that it shares nothing \omterm{def:g10:tuples-sets-comprehensions:mutable}{mutable} with the original. In Python, \code{list(a)} and \code{a[:]} make shallow copies,
and \code{copy.deepcopy(a)} makes a deep copy.
\end{definition}

\omcode{code/grade-10/02-tuples-sets-comprehensions/listings/copy_nested.py}{1}{12}{A grid of two rows, copied two
ways.}\label{lst:g10:tuples-sets-comprehensions:copy-nested}

\omoutput{code/grade-10/02-tuples-sets-comprehensions/outputs/copy_nested.txt}{The change through the \omterm{def:g10:tuples-sets-comprehensions:copy}{shallow copy}
reaches the original grid.}\label{out:g10:tuples-sets-comprehensions:copy-nested}

\begin{omfigure}
\begin{tikzpicture}[font=\small, cell/.style={draw, minimum width=6mm, minimum height=5.5mm, inner sep=0pt,
    font=\footnotesize\ttfamily}, nm/.style={draw, rounded corners=2pt, fill=omDef!8, minimum height=5.5mm,
    font=\footnotesize\ttfamily}, dot/.style={draw, minimum width=6mm, minimum height=5.5mm, inner sep=0pt}]
  % grid on the left of the shared rows, shallow on the right
  \node[nm] (g) at (0,0.6) {grid};
  \node[dot] (g1) at (1.4,0.875) {}; \node[dot, below=0pt of g1] (g2) {};
  \draw[ptr] (g) -- (g1.south west);
  \foreach \v [count=\i] in {9,0,0} \node[cell] (r\i) at (2.4+0.6*\i,1.5) {\v};
  \foreach \v [count=\i] in {0,0,0} \node[cell] (q\i) at (2.4+0.6*\i,-0.3) {\v};
  \node[dot] (s1) at (5.2,0.875) {}; \node[dot, below=0pt of s1] (s2) {};
  \node[nm] (s) at (6.7,0.6) {shallow};
  \draw[ptr] (s) -- (s1.south east);
  \fill (g1.center) circle (1.3pt); \fill (g2.center) circle (1.3pt);
  \fill (s1.center) circle (1.3pt); \fill (s2.center) circle (1.3pt);
  \draw[ptr] (g1.center) -- (r1.west); \draw[ptr] (g2.center) -- (q1.west);
  \draw[ptr] (s1.center) -- (r3.east); \draw[ptr] (s2.center) -- (q3.east);
  % deep, with rows of its own
  \node[nm] (d) at (8.6,0.6) {deep};
  \node[dot] (d1) at (9.8,0.875) {}; \node[dot, below=0pt of d1] (d2) {};
  \draw[ptr] (d) -- (d1.south west);
  \foreach \v [count=\i] in {0,0,0} \node[cell, fill=omProp!10] (e\i) at (10.8+0.6*\i,1.5) {\v};
  \foreach \v [count=\i] in {0,0,7} \node[cell, fill=omProp!10] (f\i) at (10.8+0.6*\i,-0.3) {\v};
  \fill (d1.center) circle (1.3pt); \fill (d2.center) circle (1.3pt);
  \draw[ptr] (d1.center) -- (e1.west); \draw[ptr] (d2.center) -- (f1.west);
\end{tikzpicture}

{\small After the listing: \code{grid} and its \omterm{def:g10:tuples-sets-comprehensions:copy}{shallow copy} are two outer lists sharing the same two rows, so the
\code{9} written through \code{shallow} is seen through \code{grid}; the \omterm{def:g10:tuples-sets-comprehensions:copy}{deep copy} has rows of its own.}
\end{omfigure}

\section{Sets}\label{sec:g10:tuples-sets-comprehensions:sets}

\begin{definition}[Set (Python type)]\label{def:g10:tuples-sets-comprehensions:set}
A \emph{set}\index{set (type)} is a collection of values without order and without repeats. It is written between
braces, \code{{1, 2, 3}}, or built from any sequence with \code{set(t)}; the empty set is \code{set()}, since
\code{{}} is an empty dictionary. The test \code{x in s} says whether \code{x} belongs to \code{s}, \code{s.add(x)}
adds an element, and the operators \code{&}, \code{|} and \code{-} give the intersection, the union and the
difference of two sets.
\end{definition}

\omcode{code/grade-10/02-tuples-sets-comprehensions/listings/sets.py}{1}{13}{The genres two readers like, as
\omterm{def:g10:tuples-sets-comprehensions:set}{sets}.}\label{lst:g10:tuples-sets-comprehensions:sets}

\omoutput{code/grade-10/02-tuples-sets-comprehensions/outputs/sets.txt}{The repeated \code{"mystery"} is kept once;
the order in which a \omterm{def:g10:tuples-sets-comprehensions:set}{set} prints is Python's choice.}\label{out:g10:tuples-sets-comprehensions:sets}

Two properties of \omterm{def:g10:tuples-sets-comprehensions:set}{sets} are worth keeping in mind. First, a \omterm{def:g10:tuples-sets-comprehensions:set}{set} has no positions: \code{s[0]} is an error, and the
order in which a \omterm{def:g10:tuples-sets-comprehensions:set}{set} of strings prints can differ from one computer to another. To print a \omterm{def:g10:tuples-sets-comprehensions:set}{set} in a fixed order,
print \code{sorted(s)}, which is a list. Second, the test \code{x in s} is fast, whatever the size of the \omterm{def:g10:tuples-sets-comprehensions:set}{set},
whereas \code{x in t} on a list compares \code{x} with the elements one by one; grade 11 measures the difference.

The elements of a \omterm{def:g10:tuples-sets-comprehensions:set}{set}, like the keys of a dictionary, must be \omterm{def:g10:tuples-sets-comprehensions:mutable}{immutable}: Python files each element by a number
computed from its value, and that number would be wrong if the value could change afterwards. A \omterm{def:g10:tuples-sets-comprehensions:tuple}{tuple} can be such a
key; a list cannot.

\omcode{code/grade-10/02-tuples-sets-comprehensions/listings/tuple_key.py}{1}{5}{The cells of a maze, keyed by
their coordinates.}\label{lst:g10:tuples-sets-comprehensions:tuple-key}

\omoutput{code/grade-10/02-tuples-sets-comprehensions/outputs/tuple_key.txt}{Python calls a value that
may serve as a key \emph{hashable}; a list is not.}\label{out:g10:tuples-sets-comprehensions:tuple-key}

Now the question of the introduction. The function below gathers the books one member borrowed into a \omterm{def:g10:tuples-sets-comprehensions:set}{set}, so that
a book borrowed twice counts once, and the operator \code{&} does the rest.

\omcode{code/grade-10/02-tuples-sets-comprehensions/listings/shared_books.py}{1}{20}{Which books have members 40 and
55 both borrowed?}\label{lst:g10:tuples-sets-comprehensions:shared-books}

\omoutput{code/grade-10/02-tuples-sets-comprehensions/outputs/shared_books.txt}{Member 40 is Dara, member 55 is
Umar.}\label{out:g10:tuples-sets-comprehensions:shared-books}

\begin{omfigure}
\begin{tikzpicture}[font=\small]
  \draw[thick, fill=omDef!10, fill opacity=0.7] (0,0) circle (1.5);
  \draw[thick, fill=omProp!14, fill opacity=0.7] (2.0,0) circle (1.9);
  \draw[thick] (0,0) circle (1.5);
  \node at (-0.8,0) {$1$};
  \node at (0.85,0) {$5$};
  \node at (2.9,0) {$6$};
  \node[above] at (-0.6,1.5) {Dara (6 books)};
  \node[above] at (2.6,1.9) {Umar (11 books)};
  \node[font=\footnotesize, align=left, anchor=west] at (4.3,0) {union: $1 + 5 + 6 = 12$ books\\intersection:
    $5$ books};
\end{tikzpicture}

{\small The two members' \omterm{def:g10:tuples-sets-comprehensions:set}{sets} of books, with the counts printed by the run (the five shared books are 16, 57, 62,
71 and 75). The union holds $6 + 11 - 5 = 12$
books: the shared ones would otherwise be counted twice.}
\end{omfigure}

\section{Comprehensions}\label{sec:g10:tuples-sets-comprehensions:comprehensions}

Many loops build a new collection from an old one: the squares of a list of numbers, the even ones, the first
letter of each word. Python has a shorter notation for them.

\begin{definition}[Comprehension]\label{def:g10:tuples-sets-comprehensions:comprehension}
A \emph{comprehension}\index{comprehension} is an expression that builds a list, a \omterm{def:g10:tuples-sets-comprehensions:set}{set} or a dictionary from a
sequence: \code{[e for x in t]} is the list of the values of the expression \code{e} for each element \code{x} of
\code{t}, in order; \code{[e for x in t if c]} keeps only the elements for which the condition \code{c} is true.
With braces, \code{{e for x in t}} builds a \omterm{def:g10:tuples-sets-comprehensions:set}{set} and \code{{k: v for x in t}} a dictionary.
\end{definition}

\begin{omfigure}
\begin{tikzpicture}[font=\small]
  \begin{scope}[every node/.style={font=\ttfamily, inner sep=1pt, anchor=base west}]
    \node (o) at (0,0) {[};
    \node[text=omThm] (e) at (o.base east) {n * n};
    \node (f) at ([xshift=4pt]e.base east) {for};
    \node[text=omDef] (l) at ([xshift=4pt]f.base east) {n in numbers};
    \node (i) at ([xshift=4pt]l.base east) {if};
    \node[text=omProp!80!black] (c) at ([xshift=4pt]i.base east) {n \% 2 == 0};
    \node (cl) at (c.base east) {]};
  \end{scope}
  \foreach \n/\t in {e/{3. put this\\value in}, l/{1. for each\\element}, c/{2. if the\\condition holds}}
    \draw[decorate, decoration={brace, amplitude=4pt, mirror}] ([yshift=-3pt]\n.south west) --
      ([yshift=-3pt]\n.south east) node[midway, below=5pt, align=center, font=\footnotesize] {\t};
  \node[font=\footnotesize] (eq) at ([xshift=12mm]cl.east) {$=$};
  \node[font=\footnotesize\ttfamily, align=left, anchor=west] at ([xshift=4mm]eq.east)
    {result = []\\for n in numbers:\\\ \ \ \ if n \% 2 == 0:\\\ \ \ \ \ \ \ \ result.append(n * n)};
\end{tikzpicture}

{\small A \omterm{def:g10:tuples-sets-comprehensions:comprehension}{comprehension} is read from the middle: the loop, then the condition, then the expression, which comes
first in the writing. On the right, the loop it stands for.}
\end{omfigure}

\omcode{code/grade-10/02-tuples-sets-comprehensions/listings/comprehensions.py}{1}{15}{List, \omterm{def:g10:tuples-sets-comprehensions:set}{set} and dictionary
\omterm{def:g10:tuples-sets-comprehensions:comprehension}{comprehensions}.}\label{lst:g10:tuples-sets-comprehensions:comprehensions}

\omoutput{code/grade-10/02-tuples-sets-comprehensions/outputs/comprehensions.txt}{The last \omterm{def:g10:tuples-sets-comprehensions:comprehension}{comprehension} has two
\code{for}: the second varies fastest, as in two nested
loops.}\label{out:g10:tuples-sets-comprehensions:comprehensions}

The \omterm{def:g10:tuples-sets-comprehensions:set}{set} \omterm{def:g10:tuples-sets-comprehensions:comprehension}{comprehension} on line 9 removes the repeated \code{m} and \code{g} by itself; the dictionary
\omterm{def:g10:tuples-sets-comprehensions:comprehension}{comprehension} on line 10 maps each word to its length.

\begin{example}[The books of a member, in one line]\label{ex:g10:tuples-sets-comprehensions:books-of}
The function \code{books_of} of the town listing is the \omterm{def:g10:tuples-sets-comprehensions:set}{set} \omterm{def:g10:tuples-sets-comprehensions:comprehension}{comprehension}
\code{{loan["book_id"] for loan in loans if loan["member_id"] == member}}: for each loan, if it is the member's,
put its book in the \omterm{def:g10:tuples-sets-comprehensions:set}{set}.
\end{example}

\begin{remark}[When a loop is clearer]\label{rem:g10:tuples-sets-comprehensions:loop-clearer}
A \omterm{def:g10:tuples-sets-comprehensions:comprehension}{comprehension} is for building a collection and nothing else. When each step must also print, count, stop early
(\code{break}) or update several variables, a loop is clearer; and a \omterm{def:g10:tuples-sets-comprehensions:comprehension}{comprehension} that no longer fits on one line
is usually better written as a loop. The test is the reader: if the \omterm{def:g10:tuples-sets-comprehensions:comprehension}{comprehension} cannot be read aloud as one
sentence, it should be a loop.
\end{remark}

\begin{method}[Choosing a collection]\label{met:g10:tuples-sets-comprehensions:choose}
\begin{enumerate}
  \item A fixed group of values that belong together (a point, a date, a record): a \omterm{def:g10:tuples-sets-comprehensions:tuple}{tuple} or a \omterm{def:g10:tuples-sets-comprehensions:named-tuple}{named tuple}.
  \item Values in an order that matters, possibly repeated, possibly changing: a list.
  \item Membership questions (``is it there?'', ``what is common?'') with no order and no repeats: a \omterm{def:g10:tuples-sets-comprehensions:set}{set}.
  \item Values looked up by a key: a dictionary.
\end{enumerate}
\end{method}

\section{Exercises}\label{sec:g10:tuples-sets-comprehensions:exercises}

\begin{exercise}[$\star$]\label{exo:g10:tuples-sets-comprehensions:1}
What do these lines print? \code{a, b, c = (1, 2, 3)}, then \code{a, c = c, a}, then \code{print(a, b, c)}. What
happens with \code{x, y = (1, 2, 3)}?
\end{exercise}

\begin{exercise}[$\star$]\label{exo:g10:tuples-sets-comprehensions:2}
With \code{s = {1, 2, 3, 4}} and \code{t = {3, 4, 5}}, give \code{s & t}, \code{s | t}, \code{s - t},
\code{t - s}, \code{len(s | t)} and \code{5 in s}.
\end{exercise}

\begin{exercise}[$\star$]\label{exo:g10:tuples-sets-comprehensions:3}
Write a function \code{first_last(t)} that returns the first and the last element of a non-empty list as a \omterm{def:g10:tuples-sets-comprehensions:tuple}{tuple},
with its specification, and a call that unpacks the result into two names.
\end{exercise}

\begin{exercise}[$\star$]\label{exo:g10:tuples-sets-comprehensions:4}
What does this print? \code{t = [1, 2]}, \code{u = t}, \code{v = t[:]}, \code{u.append(3)}, \code{v.append(4)},
then \code{print(t, u, v, t is u, t is v, t == u)}. Draw the memory as in the figure on \omterm{def:g10:tuples-sets-comprehensions:alias}{aliases}.
\end{exercise}

\begin{exercise}[$\star$]\label{exo:g10:tuples-sets-comprehensions:5}
Explain why \code{len(a | b) == len(a) + len(b) - len(a & b)} for any two \omterm{def:g10:tuples-sets-comprehensions:set}{sets}, and check it on the run of the
town listing. When is \code{len(a | b)} equal to \code{len(a) + len(b)}?
\end{exercise}

\begin{exercise}[$\star\star$]\label{exo:g10:tuples-sets-comprehensions:6}
A student writes \code{colours = ("red")} and then \code{for c in colours: print(c)}, hoping to print one colour.
What does the loop print, and why? Correct the first line.
\end{exercise}

\begin{exercise}[$\star\star$]\label{exo:g10:tuples-sets-comprehensions:7}
Write \code{common_books(loans, a, b)}, which returns the \omterm{def:g10:tuples-sets-comprehensions:set}{set} of the ids of the books borrowed both by member
\code{a} and by member \code{b}, using two \omterm{def:g10:tuples-sets-comprehensions:set}{set} \omterm{def:g10:tuples-sets-comprehensions:comprehension}{comprehensions}. Check it on members $40$ and $55$.
\end{exercise}

\begin{exercise}[$\star\star$]\label{exo:g10:tuples-sets-comprehensions:8}
To make a $3 \times 3$ grid of zeros, a student writes \code{grid = [[0] * 3] * 3}, then \code{grid[0][0] = 1}, and
prints \code{[[1, 0, 0], [1, 0, 0], [1, 0, 0]]}. Explain with a memory drawing, and build the grid with a
\omterm{def:g10:tuples-sets-comprehensions:comprehension}{comprehension} instead.
\end{exercise}

\begin{exercise}[$\star\star$]\label{exo:g10:tuples-sets-comprehensions:9}
What do these print? \code{{w: w.upper() for w in ["a", "b"]}}, then
\texttt{\{n \% 3 for n in range(10)\}}, then \code{{n: n * n for n in range(4) if n != 2}}.
\end{exercise}

\begin{exercise}[$\star\star$]\label{exo:g10:tuples-sets-comprehensions:10}
A \omterm{def:g10:tuples-sets-comprehensions:comprehension}{comprehension} can list all the subsets of \code{["a", "b", "c", "d"]}. Without running it, how many are there?
Explain the count by deciding, element by element, whether it is in or out, and say how many subsets have exactly
two elements.
\end{exercise}

\begin{exercise}[$\star\star$]\label{exo:g10:tuples-sets-comprehensions:11}
Write \code{distinct_authors(books)}, which returns the \omterm{def:g10:tuples-sets-comprehensions:set}{set} of the authors of a list of book records, and
\code{authors_in_genre(books, genre)}, which returns the sorted list of the authors of at least one book of the
genre. How many distinct authors does the town library have?
\end{exercise}

\begin{exercise}[$\star\star\star$]\label{exo:g10:tuples-sets-comprehensions:12}
Without \omterm{def:g10:tuples-sets-comprehensions:set}{sets}, the common books of two lists \code{a} and \code{b} are found by taking each element of \code{a} and
comparing it with the elements of \code{b} one by one until a match is found. \emph{(a)} What is the largest number
of comparisons for lists of $6$ and $11$ books? \emph{(b)} When is it reached? \emph{(c)} For two lists of $n$
elements each, what is the largest number, and how does it change when $n$ doubles?
\end{exercise}

\begin{exercise}[$\star\star\star$]\label{exo:g10:tuples-sets-comprehensions:13}
Write \code{transpose(table)}, which returns the transpose of a rectangular list of rows (row $i$ of the result is
column $i$ of the table), as a single nested \omterm{def:g10:tuples-sets-comprehensions:comprehension}{comprehension}. Check that transposing twice gives back the table.
\end{exercise}

\begin{exercise}[$\star\star\star$]\label{exo:g10:tuples-sets-comprehensions:14}
A game computes the next state of a grid of \code{0} and \code{1} by flipping every cell. Its function starts with
\code{new = list(grid)} and writes \code{new[r][c] = 1 - grid[r][c]} for every cell. The game shows the right next
state, but the previous state, kept by the caller for an ``undo'' button, has changed too. Find the cause and fix it.
\end{exercise}

\begin{exercise}[$\star\star\star$]\label{exo:g10:tuples-sets-comprehensions:15}
The \emph{Jaccard index} of two \omterm{def:g10:tuples-sets-comprehensions:set}{sets} $A$ and $B$, not both empty, is $|A \cap B| / |A \cup B|$: the share of their
union that they have in common. Write \code{jaccard(a, b)}, which returns it as an exact fraction
(\code{fractions.Fraction}), with an assertion on its \omterm{def:g10:specifications-and-tests:pre-post}{precondition}. What are its smallest and largest possible
values, and when are they reached?
\end{exercise}

\section{Problem: Readers Alike}\label{sec:g10:tuples-sets-comprehensions:problem}

\begin{problem}[{Weekend problem --- which two members read most alike?}]\label{pb:g10:tuples-sets-comprehensions:1}
The town library wants to start a reading club and pair up members with similar tastes. It measures how alike two
readers are with the Jaccard index of their \omterm{def:g10:tuples-sets-comprehensions:set}{sets} of borrowed books (\cref{exo:g10:tuples-sets-comprehensions:15}).
The program \code{exercises/readers_alike.py} of this chapter's code does the computation; the questions follow it
step by step.

\needspace{4\baselineskip}\medskip\noindent\textbf{Part I --- One set per member.}
\begin{enumerate}
  \item The program builds a dictionary \code{books} from each member id to a \omterm{def:g10:tuples-sets-comprehensions:set}{set} of book ids, with
    \code{books.setdefault(m, set()).add(b)}. Explain what \code{setdefault} does here (look it up with
    \code{help(dict.setdefault)}), and why a \omterm{def:g10:tuples-sets-comprehensions:set}{set} is used rather than a list.
  \item The library recorded $400$ loans. Run the program: how many members borrowed at least one book, and how many
    different books were borrowed?
  \item Why can the total of the sizes of the members' \omterm{def:g10:tuples-sets-comprehensions:set}{sets} be smaller than $400$?
  \item Write the dictionary \code{books} as a dictionary \omterm{def:g10:tuples-sets-comprehensions:comprehension}{comprehension} whose values are \omterm{def:g10:tuples-sets-comprehensions:set}{set} \omterm{def:g10:tuples-sets-comprehensions:comprehension}{comprehensions}, and
    say why this version is slower than the loop of the program.
\end{enumerate}

\needspace{4\baselineskip}\medskip\noindent\textbf{Part II --- Two members.}
\begin{enumerate}[resume]
  \item From the run of the town listing, give the sizes of the \omterm{def:g10:tuples-sets-comprehensions:set}{sets} of members $40$ and $55$, of their
    intersection and of their union.
  \item Compute the Jaccard index of members $40$ and $55$ as a fraction, then as a decimal to three places.
  \item Give the Jaccard index of a member with herself, and of two members with no book in common.
  \item Member $32$ borrowed one book, member $47$ two, and they share one. Compute their index.
\end{enumerate}

\needspace{4\baselineskip}\medskip\noindent\textbf{Part III --- All the pairs.}
\begin{enumerate}[resume]
  \item The program lists the pairs with the \omterm{def:g10:tuples-sets-comprehensions:comprehension}{comprehension} \code{[... for i in ids for j in ids if i < j]}. Why the
    condition \code{i < j}?
  \item How many pairs are there for $60$ members? Give the count as a formula in the number $n$ of members.
  \item The pairs are sorted with the key \code{(-value, i, j)}. Explain how this \omterm{def:g10:tuples-sets-comprehensions:tuple}{tuple} orders them, and why the two
    last components are there.
  \item Run the program. Which pair comes first among all $60$ members, and with which index?
  \item Explain why this pair says little about taste.
  \item The program then keeps only the members with at least $5$ books. How many members and how many pairs are
    left?
\end{enumerate}

\needspace{4\baselineskip}\medskip\noindent\textbf{Part IV --- The closest readers.}
\begin{enumerate}[resume]
  \item Which pair comes first among the members with at least $5$ books, and with which index?
  \item Which pairs come next, and how far behind?
  \item Show that the Jaccard index of a \omterm{def:g10:tuples-sets-comprehensions:set}{set} of $6$ elements and a \omterm{def:g10:tuples-sets-comprehensions:set}{set} of $11$ elements is at most $6/11$, and say when
    it reaches that value. How far is the pair $40$ and $55$ from it?
  \item How many comparisons of two \omterm{def:g10:tuples-sets-comprehensions:set}{sets} does the program make for the members with at least $5$ books?
  \item Check the first place with a different method: \code{itertools.combinations(ids, 2)} lists the same pairs;
    take the \code{max} of the indices.
  \item State the result: the two members whose borrowed books are most alike, with their index as a fraction and
    a decimal, and the number of books each borrowed.
\end{enumerate}
\end{problem}
